\section{Introduction}
\label{sec:introduction}

Let $\Lambda_{n,w}=(\Z/n\Z)\times(\Z/w\Z)$ be a periodic square-lattice
rectangle, and assign a spin $\sigma_x\in S^3\subset\R^4$ to each site.
We use normalized sphere measure $\dd\omega$ and the probability law
\begin{equation}
\label{eq:measure}
 \dd\mu_{\beta;n,w}(\sigma)=Z_\beta(n,w)^{-1}
 \exp\left(\beta\sum_{\langle xy\rangle}\sigma_x\cdot\sigma_y\right)
 \prod_{x\in\Lambda_{n,w}}\dd\omega(\sigma_x).
\end{equation}
Every unoriented nearest-neighbor bond is counted once. Periods are
even and at least four. Equivalently, the action, up to a constant, is
$\frac\beta2\sum_{\langle xy\rangle}|\sigma_x-\sigma_y|^2$.
The parameter $\beta>0$ is the bare inverse coupling. The number four
refers to the internal spin space; spacetime is two-dimensional.

The principal question is how rapidly the model loses memory over long
spatial separations. Three assertions must be distinguished: absence of
magnetization, exponential decay of a particular spin correlation, and a
positive gap for the full transfer operator. The third concerns every
local observable, including the rotation-invariant bond energy
$\sigma_x\cdot\sigma_y$. Its low-temperature scale is the subject of this
paper.

The absence of ordinary ordering in two dimensions is a classical
consequence of continuous symmetry. Mermin and Wagner established the
foundational nonordering result for quantum Heisenberg systems, and
Mermin gave the corresponding classical argument~\cite{MW,Mermin1967}.
McBryan and Spencer obtained quantitative algebraic upper bounds by
complex spin rotations~\cite{McBryanSpencer}; later extensions substantially
weakened the regularity assumptions on the interaction~\cite{GagnebinVelenik}.
These results leave open whether the decay is exponential. For two spin
components the distinction is essential: the low-temperature phase
predicted by Berezinskii, Kosterlitz and Thouless and established rigorously
by Fr\"ohlich and Spencer has no positive exponential spin-decay
rate~\cite{Berezinskii,KosterlitzThouless,FS}.

At sufficiently high temperature, local Ward inequalities give exponential
spin-correlation decay~\cite[Theorem~3.2]{AizenmanSimonWard}.
For non-Abelian sigma models, perturbative renormalization further
predicts asymptotic freedom and mass generation~\cite{Polyakov,BrezinZinnJustin,BrezinZinnJustinLeGuillou}.
In the normalization \eqref{eq:measure}, the predicted $O(4)$ inverse
length is $\sqrt\beta e^{-\pi\beta}$. Bethe-Ansatz and $S$-matrix
calculations of the exact continuum mass for $O(3)$ and $O(4)$, and
subsequently general $O(n)$, provide a
complementary description~\cite{HasenfratzMaggioreNiedermayer,HasenfratzNiedermayer}.
They concern a continuum theory and its renormalization scale; they do
not by themselves prove a gap for the specified lattice model or its
convergence to that theory. The all-temperature lattice decay question
is recorded as open in the 2025 account of Aru, Garban and
Sep\'ulveda~\cite[Section~1]{AGS}. We address the $O(4)$ model with its full
local-observable transfer spectrum and keep the continuum issue separate.

OpenAI~\cite[Theorem~1.1 and Corollary~1.2]{OpenAIClassicalONDecay}
proves exponential spin-correlation decay for each fixed integer internal
spin dimension at least three and each finite $\beta>0$. The bounds are
uniform over finite free-boundary nearest-neighbor square-lattice subgraphs
with edge strengths in $[0,\beta]$ and pass to subsequential local weak
limits of free boxes with constant strength $\beta$. They neither identify
these limits with the periodic state used here nor assert a full
local-observable transfer gap.

The reflected Hilbert-space construction originates in the work of
Osterwalder and Schrader~\cite{OS73,OS75}; the elementary discrete
construction needed here is given in Section~\ref{pre:part-12-2}.
Reflection positivity associates to a periodic thermodynamic limit a
Hilbert space, a vacuum vector $\Omega$, and a positive one-step
translation contraction $T_\beta$, with $T_\beta\Omega=\Omega$.
The Hilbert space is the completion of local observables on a time
half-space under their reflected two-point form. Define the full mass
gap by
\begin{equation}
\label{eq:mass-definition}
 \gap(\beta)=-\log\norm{T_\beta\big|_{\Omega^\perp}}.
\end{equation}
This definition includes every sector generated by local observables,
including invariant observables such as bond energies. A spin
two-point estimate alone would not prove a lower bound for this gap.
We abbreviate the natural inverse length by
$\rho(\beta)=\sqrt\beta e^{-\pi\beta}$.

The theorem also gives bounds in terms of a coarse lattice. The block
transformation used below introduces a unit spin for each $L\times L$
block through a normalized probability kernel, and then integrates the
finer spins exactly. After $N$ steps, one coarse lattice spacing
represents $L^N$ microscopic spacings. The negative logarithm of the
resulting density is its effective action. The \emph{kinetic coupling}
is the bulk coefficient extracted from the quadratic cost of a slowly
varying spin configuration in that action, in the same inverse-coupling convention
as $\beta$. At the initial step it is $\beta$. The
\emph{terminal coupling} is its value when blocking stops.

\begin{theorem}
\label{thm:main}
There are finite constants $\beta_0,c,C$, with $0<c<C$, such that for
every $\beta\ge\beta_0$ the periodic square measures have a unique local
limit, and its full transfer gap satisfies
\begin{equation}
\label{eq:target}
 c\sqrt\beta\,e^{-\pi\beta}\le\gap(\beta)
            \le C\sqrt\beta\,e^{-\pi\beta}.
\end{equation}
There are also a fixed dyadic blocking factor $L$, a fixed large
terminal coupling $H$, and finite integers $N_0,M_0$ with the following
property. For every $N\ge N_0$ one can choose a bare coupling
$\beta_N$ whose exact $N$-step block transformation has terminal
kinetic coupling $H$, and
\begin{equation}
\label{eq:depth-bounds}
 \frac{\log2}{M_0L^N}\le\gap(\beta_N)
                  \le\frac{\log8}{L^N}.
\end{equation}
The same bounds hold uniformly for terminal couplings in a fixed small
interval about $H$, after increasing $M_0$ and $N_0$ if necessary.
\end{theorem}

\begin{corollary}[Full gap at every positive temperature]
\label{cor:all-temperature}
For every finite $\beta>0$, the measures \eqref{eq:measure} on periodic
squares have a unique local limit. There are absolute constants $c,C>0$
such that its full transfer gap obeys
\[
 \gap(\beta)\ge c\exp\{-\pi\beta-
 C\sqrt{(1+\beta)\log(2+\beta)}\}>0.
\]
\end{corollary}

Corollary~\ref{cor:all-temperature} follows from the preliminary
all-observable mixing estimate in Section~\ref{pre:section}, without
using the sharp-scale comparison. Uniqueness here refers to the periodic
local limit; no assertion about every nonperiodic Gibbs state is included.
For the $O(4)$ periodic state this gives a positive resolution of the
all-temperature lattice mass-generation conjecture discussed above.
Theorem~\ref{thm:main} gives two-sided order bounds, not an exact
asymptotic prefactor.

In particular, fix a positive physical reference length
$\ell_{\mathrm{ref}}$ and declare that $L^N$ microscopic steps have
this length. The lattice spacing is $a_N=\ell_{\mathrm{ref}}L^{-N}$.
Equation~\eqref{eq:depth-bounds} gives
\begin{equation}
\label{eq:physical-main}
 \frac{\log2}{M_0\ell_{\mathrm{ref}}}
 \le \frac{\gap(\beta_N)}{a_N}
 \le \frac{\log8}{\ell_{\mathrm{ref}}}.
\end{equation}
The units have been fixed by the block scale and terminal coupling,
independently of the mass. Up to bounded positive multiplicative
factors, the two-loop calculation identifies this choice with the
natural asymptotically free scale for
\eqref{eq:measure}. These are bounds for the cutoff theories; neither
convergence of their masses nor existence of a continuum theory is
assumed.

\subsection{Methods and their relation to earlier work}

The proof joins three traditions. Local rotation identities turn
continuous symmetry into estimates on current fluctuations, as in the
classical Ward-identity approach~\cite{Mermin1967,DriesslerLandauPerez,AizenmanSimonWard}.
Here the quaternion commutator supplies a quantitative decrease of
stiffness across scales. The later passage from charged observables to
invariant ones uses Ginibre's formulation of the Griffiths inequalities
for Ising and plane-rotator systems, the FKG association mechanism, and
embedded reflection variables
\cite{Ginibre1970,FortuinKasteleynGinibre1971,Wolff1989}, followed by two
successive conditioning arguments. We prove the finite forms needed
below and retain the covariances of conditional means at both stages.

Reflection positivity connects Euclidean lattice correlations with a
positive transfer operator. Classical lattice reflection and chessboard
methods were developed systematically by Fr\"ohlich, Israel, Lieb and
Simon~\cite{FILS78,FILS80}. The finite-volume argument uses the complete
trace, so multiplicities and invariant sectors enter automatically.
Its role differs from that of an infrared bound: it turns a small
partition-function defect on one box into a spectral estimate uniform
in larger circumferences. The positivity calculations and the limiting
Hilbert-space argument are included in Section~\ref{sec:finite-volume}.

The block-variable viewpoint goes back to Kadanoff and
Wilson~\cite{Kadanoff,Wilson}. Constructive results for sigma models have
required careful distinctions of regime and model. Kupiainen established
a large-$n$ expansion and a mass gap for sufficiently large $n$ above the
spherical model's critical temperature~\cite{Kupiainen}; Gaw\k{e}dzki and
Kupiainen constructed a finite-volume continuum limit for a hierarchical
model~\cite{GawedzkiKupiainen}; Mitter and
Ramadas constructed the Wilson trajectory in perturbation
 theory~\cite{MitterRamadas}. Ba\l aban developed small-field
renormalization and variational-background methods for vector spins
\cite{BalabanFlow,BalabanVariational}. For $SU(2)$, identified with $S^3$,
Dybalski, Stottmeister and Tanimoto proved a small-field
variational result for one block step, uniform in the box size~\cite{DybalskiStottmeisterTanimotoVariational}.
These results are relevant predecessors of the block analysis; their
hypotheses and conclusions do not supply the exact iterated $O(4)$
density estimates required here.

Normalized stochastic block transformations for sigma models also appear
in the perfect-action construction of Hasenfratz and Niedermayer
\cite[Section~2.1, Eqs.~(3)--(5)]{HasenfratzNiedermayerPerfect1994}.
The observation kernel used here differs, and its iterated estimates
are proved below. The renormalization step is an exact integration
through that probability kernel. Its estimates concern the complete density, including
rough configurations, rather than only a formal expansion near aligned
spins. Perturbative coefficients determine the scale only after these
uniform estimates are available. The comparison of two cutoffs then
requires equal terminal kinetic couplings. Since the retained expansion
has signed terms, we compare complete positive densities on a common
set of large probability under both laws. These are the two interfaces
that connect the coefficient calculation with the nonperturbative gap.

\subsection{Proof strategy}

There are two distinct changes of scale in the proof. At fixed $\beta$,
we enlarge a periodic box to control the infinite-volume spectrum.
To compare ultraviolet cutoffs, we instead change the bare coupling
and the number of blocking steps while holding the coarse coupling
fixed. The argument succeeds because the second operation preserves a
finite-box estimate to which the first operation applies.

\subsubsection{A finite box controls the full spectrum}

On a cylinder of spatial circumference $w$, let $K_w$ be the positive
transfer operator that advances time by one lattice step. If its
eigenvalues are $\lambda_1>\lambda_2\ge\cdots\ge0$, closing time after
$n$ steps gives $Z_\beta(n,w)=\operatorname{Tr}K_w^n$. The ratio
\[
 P_\beta(n,w)=\frac{Z_\beta(2n,w)}{Z_\beta(n,w)^2}
 =\sum_j\left(\frac{\lambda_j^n}{\sum_i\lambda_i^n}\right)^2
\]
is the purity of the normalized spectral weights. When it is close to
one, almost all the trace weight belongs to the largest eigenvalue.
This controls the sum of the excited-state contributions, including
their multiplicities. Writing
\[
 s_n(w)=\sum_{j\ge2}(\lambda_j/\lambda_1)^n,
 \qquad s_{2n}(w)\le s_n(w)^2,
\]
shows the elementary source of the improvement under time doubling.

The remaining issue is that enlarging the spatial circumference could
introduce new low-energy states. Transfer positivity in both coordinate
directions resolves this issue through exact identities among
rectangular partition functions. For a square define
\begin{equation}
 \Delta_\beta(n)=4\log Z_\beta(n,n)-\log Z_\beta(2n,2n).
 \label{intro:Delta}
\end{equation}
By interchange of the coordinate axes,
\[
 e^{-\Delta_\beta(n)}=P_\beta(n,n)^2P_\beta(n,2n).
\]
The second purity is evaluated on the already-doubled rectangle.
A product close to one therefore controls both directions.
Section~\ref{sec:finite-volume} proves
\[
 \Delta_\beta(2n)\le64\Delta_\beta(n)^2,
 \qquad
 \Delta_\beta(n)\le\eta_*:=2^{-10}
 \ \Longrightarrow\ \gap(\beta)\ge\frac{\log2}{n}.
\]
The recurrence applies in the displayed small-defect regime.
Repeated squaring gives decay exponential in the increasing side
length. It controls the full transfer spectrum, including every
symmetry sector, as the circumference tends to infinity.
This is the box-doubling part of the proof.

The same quantity is useful for comparing cutoffs. Any contribution
$f_\beta nw$ to $\log Z_\beta(n,w)$ cancels from
\eqref{intro:Delta}. Thus the bulk free energy generated by exact
integration does not enter the finite-volume test. Our lower-bound
problem is now to make this one test small on a box whose physical
size stays bounded as $\beta$ increases.

\subsubsection{The preliminary estimate: currents, stiffness, and mixing}

The first source of decay is specific to the spin geometry. Identify
$S^3$ with the unit quaternions. The first-order change of each bond
energy under a spin rotation defines its current in that rotation
direction. The corresponding \emph{stiffness}
is the second-order response of the free energy to the imposed twist.
Its defining identity has the form ``direct quadratic cost minus
current covariance.'' Conditioning on the boundaries of smaller
squares gives the same structure: the stiffness of the larger square
is an average of cell stiffnesses minus a covariance. A lower bound
on these current fluctuations therefore improves an upper bound on
stiffness. No positive lower bound on the local stiffness is needed.

Rotations about different quaternion axes do not commute. Their
integration-by-parts identities consequently contain a term that
rotates the current itself. Applying these identities over a range of
wavelengths gives a stiffness reduction with leading term
$(\log \ell)/\pi$ when the side of a square is increased by a factor
$\ell$. The logarithm comes from a reciprocal-square frequency sum in
two dimensions. Starting from stiffness at most $\beta$, one can
therefore reach a small upper bound after a total logarithmic length
close to $\pi\beta$. Choosing the intermediate scale increments to
control the accumulated errors gives the length bound
\begin{equation}
 \Xi(\beta)=\exp\!\left[\pi\beta+
 C\sqrt{(1+\beta)\log(2+\beta)}\right].
 \label{intro:Xi}
\end{equation}
The extra term is a loss in this estimate, not a prediction for the
physical correlation length.

Small stiffness makes suitable local rotations inexpensive with high
probability. Coarse orientation variables then admit a coupling that
erases boundary disagreements except at rare locations. A disagreement
can travel far only through a long connected chain of such failures.
This gives decay first for observables with zero average under a common
spin rotation.

To include invariant observables, write each spin as
$(\cos\alpha_x\,z_x,\sin\alpha_x\,w_x)$, with
$\alpha_x\in[0,\pi/2]$ and $z_x,w_x\in S^1$. Given the angles, the
two circle-valued fields are independent ferromagnets. Further
conditioning on component magnitudes reduces their sign variables to
ferromagnetic Ising models. Positive-correlation inequalities relate
their general covariances to two-point functions. At both conditioning
stages the covariance of the conditional means must also be bounded;
cutting distant bonds makes those means local, with a controlled error.
This proves mixing for all local Lipschitz observables in the unpinned
ferromagnetic systems, uniformly when bonds are weakened or removed.

That uniformity also compares the mean bond energy on the $n$ and
$2n$ tori. Integrating the difference with respect to the coupling
gives the precise output of Section~\ref{pre:section}:
\begin{equation}
 0\le\Delta_\beta(n)
 \le C(1+\beta)^p n^p e^{-cn/\Xi(\beta)},
 \qquad n\ge C\Xi(\beta).
 \label{intro:preliminary}
\end{equation}
This establishes a gap at each fixed coupling. Its length bound is too
large to give a positive lower bound in the physical units of
\eqref{eq:physical-main}: the excess over
$\rho(\beta)^{-1}=\beta^{-1/2}e^{\pi\beta}$ grows without bound.
The next step uses this estimate at one reference cutoff and preserves
its conclusion at all finer cutoffs.

\subsubsection{Comparing cutoffs at a fixed coarse coupling}

Fix the blocking factor $L$ and a sufficiently large terminal inverse
coupling $H$. Every step of the exact block transformation remains in
the weak-coupling regime. The coarse law contains interactions beyond
the original nearest-neighbor energy. These are retained, together
with the contributions from configurations having large spin
differences. The comparison concerns the complete coarse laws
produced by the original lattice measures.

The basic contraction comes from averaging. In the linear prediction
of the fine field from a slowly varying coarse field, each spin
difference gains a factor $L^{-1}$, while a block contains $L^2$ sites.
This contracts the dependence on inherited higher-order interactions,
after their lower-order contributions have been extracted. Integrating
the fluctuations also generates new interactions; the estimates keep
them within the same controlled class. A quadratic
term needs a further subtraction: its response to a uniform gradient
is included in the running kinetic coupling. The remaining quadratic
terms have a cancellation that makes them contract as well. A bulk
constant is tracked separately. This is why the kinetic coupling is
the parameter that must be matched.

Let $I_H$ be a fixed small interval about $H$. A trajectory is
\emph{admitted} if every exact blocking step stays in the controlled
class of densities and coupling ranges established in
Section~\ref{rg:section}. Before constructing such trajectories,
Section~\ref{cal:section} determines the two kinetic coefficients.
The relation between depth and bare coupling is then determined, rather
than imposed, by the kinetic drift. Computing the same finite-volume
partition function directly and after exact blocking identifies the
first two drift coefficients. Together with the remainder estimates
for the exact transformation, this gives, with chronological step index $h$
and coupling $\widetilde b_h$,
\[
 \widetilde b_{h+1}=\widetilde b_h-\gamma-\frac{\gamma}{2\pi\widetilde b_h}+r_h,
 \qquad \gamma=\frac{\log L}{\pi},
\]
where, along an admitted nearest-neighbor trajectory, the total absolute
error $\sum_h|r_h|$ is bounded uniformly in the depth. Summation,
with $\widetilde b_0=\beta$ and $\widetilde b_N\in I_H$ fixed, yields
\begin{equation}
 N\log L=\pi\beta-\tfrac12\log\beta+O_{L,H}(1).
 \label{intro:calibration}
\end{equation}
The inverse-coupling correction is essential: it supplies the factor
$\sqrt\beta$ in the inverse length. Control of the leading exponential
$e^{-\pi\beta}$ alone would leave an unbounded multiplicative
uncertainty in physical units.

Section~\ref{sec:trajectories} next shows that for every terminal value
$h\in I_H$ one can choose a bare coupling $\beta_N(h)$ so that all $N$ exact steps satisfy the
required error bounds and the last coupling is $h$. Existence of these
choices is part of the construction; uniqueness is unnecessary. Compare depths $N\ge K$.
At the beginning of their last $K$ steps, the remaining interaction
data can differ by a bounded amount. Their kinetic couplings agree at
the opposite end. Forward contraction of the interaction differences,
combined with this terminal condition on the coupling, bounds the
differences by $C_He^{-c_HK}$ in the local norms controlling the
retained interactions and errors. The estimate is uniform in
$N-K$ and allows arbitrarily separated bare couplings.

Passing from local interaction bounds to partition functions requires
an additional argument. The expansion uses signed terms, so a small
absolute error in its coefficients need not be a small relative error
in its sum. Section~\ref{cmp:section} instead compares the complete
positive densities on a common set whose complement has small
probability under \emph{both} laws. Local modifications of rough
regions and an exact expansion provide this relative comparison.
The probability bounds under both laws then control the normalization
constants as well.

The errors can accumulate over the $M^2$ coarse cells in a box. This
explains both the strength and the volume restriction of the resulting
comparison. Define
\begin{equation}
 D_N(h;M)=\Delta_{\beta_N(h)}(ML^N).
 \label{intro:DN}
\end{equation}
For all sufficiently large $K$, all $N\ge K$, and all dyadic
coarse periods $M\ge M_{\min}(L,H)$, the estimate is
\begin{equation}
 |D_N(h;M)-D_K(h;M)|\le C_He^{-d_HK},
 \qquad (2M)^2\le e^{c_HK}.
 \label{intro:comparison}
\end{equation}
The constants are uniform in $h\in I_H$. The factor $2M$ occurs because
the test contains the doubled box as well. The common terminal
coupling fixes the physical comparison: the two microscopic periods
$ML^N$ and $ML^K$ both represent the same $M\times M$ coarse torus.

\subsubsection{Choose the reference box once, then refine the cutoff}

We can now see why the two estimates fit together.
Equation~\eqref{intro:calibration} gives
\[
 \frac{\Xi(\beta_K(h))}{L^K}
 =\exp\!\bigl[O_{L,H}(\sqrt{K\log K})\bigr].
\]
Thus the coarse side required by the preliminary decay bound grows
subexponentially in $K$. The comparison
\eqref{intro:comparison}, however, permits coarse sides growing
exponentially in $K$. There is room to choose a box large enough for
the first estimate and small enough for the second.

Concretely, choose dyadic $M_K$ with
\[
 \log M_K=K^{3/4}+O(1).
\]
Then
\[
 \log\frac{M_KL^K}{\Xi(\beta_K(h))}
 =K^{3/4}-O_{L,H}(\sqrt{K\log K}),
 \qquad 2\log(2M_K)=o(K).
\]
The first relation makes the decay in \eqref{intro:preliminary}
overwhelm its polynomial prefactor, so $D_K(h;M_K)\to0$ uniformly in
$h$. The second keeps both boxes inside the comparison range. The
exponent $3/4$ merely lies between $1/2$ and $1$; no significance is
attached to its particular value.

Choose one finite $K_0$ large enough that
\[
 \sup_{h\in I_H}D_{K_0}(h;M_{K_0})+C_He^{-d_HK_0}\le\eta_*.
\]
Now fix $M_0=M_{K_0}$. For every finer cutoff $N\ge K_0$,
\[
 D_N(h;M_0)\le\eta_*.
\]
Box doubling therefore gives
$\gap(\beta_N(h))\ge(\log2)/(M_0L^N)$.
The growing family $M_K$ was used only to find a single reference box.
After that choice, $M_0$ is constant, and the microscopic box
$M_0L^N$ always has physical side $M_0\ell_{\mathrm{ref}}$.
This order of choices removes the subexponential loss from the
physical mass bound. It also explains why the blocking construction
can stop at a large $H$: long-distance decay is supplied by the
preliminary argument on the reference lattice.

\subsubsection{A block correlation bounds the mass above}

The upper bound uses a different consequence of stopping at a large
$H$. Nearby terminal spins remain aligned on average, by the one-law
rarity estimate of Lemma~\ref{cmp:rarity}. For terminal
sites $0$ and $3e_1$, Lemma~\ref{lem:alignment} gives
$\mathbb E[V_0^{(1)}V_{3e_1}^{(1)}]\ge1/8$, where the superscript
denotes one spin component. Pull this observable back to the original
lattice by setting
\[
 F_X(\sigma)=\mathbb E[V_X^{(1)}\mid\sigma].
\]
Each $F_X$ is centered, bounded by one, and supported in its block of
side $L^N$. The auxiliary variables in disjoint blocks are independent
given $\sigma$, so the two $F$'s retain the same correlation. Their
supports are separated by at least $L^N$ time steps. The full transfer
gap bounds that correlation by $e^{-\gap(\beta_N(h))L^N}$, giving
the upper bound in \eqref{eq:depth-bounds}.

The proof thus gives a uniform exponential decay rate for all connected
correlations of bounded local observables in physical units, together
with a nonzero bounded correlation at a fixed physical separation.
These give the two sides of \eqref{eq:physical-main}. The admission
argument assigns every sufficiently large bare coupling a trajectory
ending in $I_H$, so the depth bounds imply the assertion for every
sufficiently large $\beta$. Section~\ref{sec:physical} states separately
the additional convergence assumptions needed to interpret the bounds
as statements about an existing continuum theory.

\paragraph{Conventions.}
The normalizations used throughout are collected here to distinguish the
physical spacing from the observation precision.
\begin{center}
\begin{tabular}{@{}ll@{}}
\toprule
Quantity & Convention\\
\midrule
Spin and measure & $\sigma_x=q_x\in S^3$; $\int_{S^3}d\omega=1$\\
Bond sum & Each unoriented nearest-neighbor bond counted once\\
Bare parameter & $\beta>0$ is inverse coupling (inverse temperature)\\
Running parameter & $b$ is the extracted kinetic inverse coupling\\
Preliminary distance & Maximum-norm lattice distance between supports\\
Transfer distance & Number of time edges between the actual supports\\
Physical spacing & $a_N=\ell_{\mathrm{ref}}L^{-N}$\\
Observation precision & $a=1$, a dimensionless kernel parameter\\
\bottomrule
\end{tabular}
\end{center}
For quaternion generators, an imaginary unit $u$ acts by
$q_x\mapsto e^{tu}q_x$ and $[u,v]=2u\times v$; the bond current
normalization is derived in Section~\ref{pre:local}.

\paragraph{Organization.}
Section~\ref{sec:finite-volume} proves the box-doubling criterion and
Section~\ref{pre:section} proves the preliminary estimate.
Sections~\ref{free:section} and~\ref{rg:section} construct and control
the exact block transformation. Section~\ref{cal:section} calibrates
its kinetic drift. Section~\ref{sec:trajectories} constructs admitted trajectories
and matches their terminal couplings; Section~\ref{cmp:section} compares the resulting
positive densities. Sections~\ref{sec:assembly}--\ref{sec:physical}
assemble the lower bound, prove the upper bound, and identify the
physical units. The regularity and support conventions required by the expansions are
specified at their points of use. The appendices are parts of the proof:
Appendix~\ref{app:analytic-block} supplies the uniform analytic estimates
for the linear block map; Appendix~\ref{supp:canonical-fixed-norm} fixes
the canonical coefficient norm and proves its diagonal contraction;
Appendix~\ref{app:rg-geometry} constructs the local factors and their
complete supports; and Appendix~\ref{supp:joint-majorants} proves the
joint Gaussian product bound and verifies it for all eight prepared-factor
classes.


